%
% Capítulo 1
%
\chapter{Introdução} \label{cap:intro}

A organização de horários entre professores e alunos é uma tarefa que, em muitos contextos, continua a ser realizada
de forma manual, recorrendo a trocas sucessivas de mensagens, folhas de cálculo ou registos informais.
Este processo tende a ser demorado, pouco flexível e propenso a erros, especialmente quando existem múltiplos alunos,
disponibilidades distintas e restrições específicas que precisam de ser respeitadas. A dificuldade aumenta quando se
pretende encontrar uma distribuição equilibrada de sessões, compatível com a disponibilidade do professor e com as
preferências ou limitações dos alunos.

A motivação para este projeto surgiu de uma experiência concreta das autoras no contexto de um centro de estudos, onde
ambas desempenham funções de explicadoras. Neste contexto, a gestão dos horários dos alunos era realizada manualmente e
exigia reajustes frequentes sempre que entrava um novo aluno ou surgiam alterações de disponibilidade. Na prática, isso
implicava refazer sucessivamente os horários já definidos, tentar encontrar novos encaixes e, muitas vezes, lidar com a
dificuldade de conciliar todos os alunos de forma equilibrada e sustentável. Esta experiência permitiu identificar de
forma clara a necessidade de uma solução que apoiasse este processo de forma mais eficiente e estruturada.

Neste contexto, surgiu a proposta do projeto \textit{Sistema de Planeamento de Horários}, cujo objetivo consiste no
desenvolvimento de uma solução capaz de apoiar a gestão e criação automática de horários entre professores e alunos.
A solução permite recolher disponibilidades, definir restrições de agendamento e criar propostas de horário de forma
estruturada, reduzindo o esforço manual associado a este processo e tornando a gestão de sessões mais clara e eficiente.

%
% Secção 1.1
%
\section{Motivação e Objetivos} \label{sec11}

A motivação principal para o desenvolvimento deste projeto prende-se com a necessidade de criar uma ferramenta capaz de
simplificar a coordenação de horários em cenários onde um professor trabalha com vários alunos e necessita de conciliar
diferentes disponibilidades ao longo da semana. Em vez de depender exclusivamente de planeamento manual, a solução proposta
procura automatizar parte desse trabalho, permitindo que os dados relevantes sejam recolhidos de forma consistente e utilizados
para construir propostas de horário com base em regras bem definidas.

O objetivo central do projeto consiste, assim, em disponibilizar um sistema que permita:
\begin{itemize}
    \item registar professores e alunos;
    \item associar alunos a um professor;
    \item recolher as disponibilidades semanais de professores e alunos;
    \item definir restrições relevantes para o agendamento;
    \item criar automaticamente propostas de horário;
    \item apresentar o horário resultante de forma legível, organizada, dinamica e facilitada.
\end{itemize}

Numa perspetiva mais ampla, o projeto pretende também servir como base para a evolução de funcionalidades futuras, como integração
com sistemas externos, persistência remota mais completa e eventual suporte a extração automática de informação a partir de documentos ou imagens.

%
% Secção 1.2
%
\section{Visão Geral da Solução} \label{sec12}

A solução desenvolvida assenta numa arquitetura composta por uma aplicação Android, responsável pela interação com o utilizador,
e por um backend com API REST, responsável pela persistência e disponibilização de dados centrais do sistema. No frontend foi utilizada
a linguagem Kotlin com Jetpack Compose para a construção da interface, enquanto no backend foi adotado Ktor para a implementação da API,
em conjunto com PostgreSQL para persistência de dados.

A aplicação distingue dois tipos principais de utilizador através de uma classe "Owner Type" desenvolvida por nós: professor e aluno.
Cada um destes perfis tem acesso a um conjunto próprio de funcionalidades.
O professor pode definir a sua disponibilidade, configurar restrições associadas ao planeamento e consultar os alunos que lhe estão associados.
O aluno pode escolher um professor e definir a sua própria disponibilidade.
A partir destes dados, o sistema pode criar uma proposta de horário semanal compatível com as disponibilidades registadas e com as restrições configuradas.
Por sua vez o professor irá avaliar o resultado final e, caso seja necessário, poderá requesitar uma nova proposta de horário.
Ao ser aceite a proposta de horário, o sistema irá registálo no Google Calendar do professor, utilizando a API do Google Calendar, e irá disponibilizar o horário especifico de cada aluno no seu email pessoal.

Na versão beta do projeto, foi dada especial atenção à validação dos fluxos principais do sistema, à integração entre frontend e backend, e à implementação
de uma primeira abordagem para a criação automática de horários. Esta fase permite demonstrar a viabilidade técnica da solução e serve como base para melhorias e refinamentos futuros.

%
% Secção 1.3
%
\section{Organização do Documento} \label{sec13}

O presente relatório encontra-se organizado da seguinte forma. No Capítulo~\ref{cap:enquadramento} é apresentado o enquadramento do problema, a motivação do projeto e o contexto em que a solução se insere.
No Capítulo~\ref{cap:arquitetura} é descrita a organização global do sistema, incluindo os principais componentes e a sua interação.
Seguidamente, são apresentados os detalhes de implementação mais relevantes, tanto ao nível do frontend como do backend, bem como os principais aspetos da API e da persistência de dados.
Posteriormente, é descrita a estratégia de testes e validação aplicada ao sistema.
Por fim, é feito um balanço do estado atual da versão beta, identificando as funcionalidades já concluídas, as limitações atuais e os próximos passos previstos para a evolução do projeto.
